\documentclass[11pt,a4paper]{article}
\usepackage[margin=24mm]{geometry}
\usepackage{fontspec}
\setmainfont{DejaVu Serif}[ItalicFont={DejaVu Serif},ItalicFeatures={FakeSlant=0.18},BoldItalicFont={DejaVu Serif Bold},BoldItalicFeatures={FakeSlant=0.18}]
\setsansfont{DejaVu Sans}
\setmonofont{DejaVu Sans Mono}[Scale=0.8]
\usepackage{babel}
\babelprovide[main,import]{ukrainian}
\usepackage{amsmath,amssymb,booktabs,longtable,graphicx,float}
\usepackage[unicode,hidelinks]{hyperref}
\usepackage{microtype}
\setlength{\emergencystretch}{3em}
\title{Як підготувати реальну контактну мережу\\для двомасштабного методу}
\author{Робоче пояснення до основної SMPT-статті}
\date{16 жовтня 2026 року}
\begin{document}
\maketitle
\noindent Цей текст пояснює перехід від відкритих записів контактів до входів методу: локальних графів, повної контрольної мережі та матриці міжспільнотних ваг. Це розбір для автора, без анотації, висновків і журнального оформлення. Відокремлюємо факти з даних, рішення підготовки та питання для майбутніх симуляцій.

\section{Що саме потрібно методу з головної статті}
Джерело вимог --- поточний англійський \texttt{main-SMPT.tex}, розділи 3--5. Стаття досліджує похибку двомасштабного наближення відносно повного CTMC, зокрема час приходу інфекції та допустимість маршрутів. Її синтетичні мережі дозволяють контролювати структуру, але для реальних даних спочатку потрібно визначити вузли, спільноти, ребра, одиниці ваг і часу.

Нехай $V=\bigsqcup_{i=0}^{M-1}V_i$, $N_i=|V_i|$. Повна модель отримує всі ребра, двомасштабна --- індуковані графи $G_i=G[V_i]$ та для $i\ne j$ лише
\[
W_{ij}=\sum_{u\in V_i,\,v\in V_j:\{u,v\}\in E}w_{uv},\qquad W_{ii}=0.
\]
Для неорієнтованого контакту вага входить у $W_{ij}$ та $W_{ji}$: це два можливі напрямки передачі. Подвоювати вагу в одному елементі не потрібно. Макроінтенсивність дорівнює
\[
\lambda_{ij}=\beta_{\mathrm{out}}W_{ij}\frac{I_i}{N_i}\frac{S_j}{N_j}.
\]
Тому $W$ є сумою ваг, а не середнім, щільністю чи сумою, вже поділеною на $N_iN_j$. Нормування повторно зменшило б інтенсивності. Розміри спільнот можуть відрізнятися; у формулі потрібні фактичні $N_i,N_j$.

Ідеальний агрегований процес відповідає іншій мережі: для кожної зв'язаної пари спільнот усі міжспільнотні пари вузлів отримують вагу $W_{ij}/(N_iN_j)$. Вона зберігає сумарну вагу, але змінює її розташування. Саме цю відмінність потрібно оцінювати на підготовленій емпіричній мережі.

\section{Чому вибрано ці дані}
Вибрано сенсорні контакти SocioPatterns High School 2013: дев'ять класів, п'ять днів спостереження, записи з кроком 20 секунд. Клас задає спільноту незалежно від моделювання. Тому не потрібно добирати кластеризацію за бажаною точністю. Дружба та щоденники з тієї самої сторінки мають інше значення; тут їх не змішуємо з фізичними контактами.

Первинна сторінка: \href{https://sociopatterns.org/datasets/high-school-contact-and-friendship-networks/}{SocioPatterns: High School 2013}.

Публікація: R.~Mastrandrea, J.~Fournet, A.~Barrat (2015), \emph{Contact Patterns in a High School: A Comparison between Data Collected Using Wearable Sensors, Contact Diaries and Friendship Surveys}, \url{https://doi.org/10.1371/journal.pone.0136497}.

На сторінці даних вказано CC BY-NC-SA, посилання веде на версію 3.0. Ці умови зберігаємо окремо від ліцензії журнальної публікації. При використанні потрібні посилання на SocioPatterns і цитування джерела. Сенсори не охоплюють контакти з особами поза вибіркою. В оригінальному дослідженні сенсори носили 327 із 379 учнів відповідних класів; отже, це спостережена підмережа, а не повний опис школи.

У \path{raw/} збережено два оригінальні файли: стиснені контакти та метадані. Їх URL, час отримання й SHA-256 зафіксовано в \path{sources.json}. Головну статтю не редагуємо; її хеш також зафіксований. Наступні розділи додаються після перевірки й перетворення самих файлів.


\section{Що показала перевірка сирих файлів}
Одиниця сирих даних --- пара людей у двадцятисекундному інтервалі. Рядок $(t,u,v,C_u,C_v)$ означає активний контакт у $[t-20,t]$, а не окрему завершену зустріч. Послідовні рядки однієї пари можуть належати одній зустрічі; складати їх тривалості можна, називати кожний рядок новою зустріччю --- ні.

У файлі 188\,508 таких записів, 327 різних учасників і 5818 різних неорієнтованих пар. Повторів ключа $(t,\min(u,v),\max(u,v))$, петель та учасників без метаданих не знайдено. Сумарна парна тривалість становить 3\,770\,160 секунд: якщо одночасно спостерігаються багато пар, кожна додає свою тривалість. Тому ця сума може перевищувати календарну тривалість спостереження.

\paragraph{329 записів метаданих проти 327 контактних вузлів.}
Ідентифікатори 2 та 478 присутні лише в метаданих. Основний граф містить 327 людей із хоча б одним контактом. Це явне визначення досліджуваної сукупності, а не твердження, що двох інших людей не існувало чи вони не контактували в реальності. У \path{prepared/nodes.csv} збережено всі 329 ID і прапорець \texttt{observed}. Варіант \texttt{population="metadata"} додає цих двох людей як ізольовані вузли. Повний CTMC не може інфікувати їх через спостережені контакти, тоді як рівномірний макроімпорт може. Тому варіанти 327 і 329 треба порівнювати окремо.

\paragraph{Неузгоджений клас учасника 28.}
Окремий файл метаданих відносить його до \texttt{2BIO2}, але всі 1321 записи, де він є кінцем контакту, містять \texttt{2BIO1}. Джерела не дозволяють встановити, яке позначення правильне. Для основної версії пріоритет має окрема таблиця метаданих: це послідовне правило, а не виправлення встановленого факту. Друга версія бере клас із контактного файлу. Жодні контакти при цьому не вилучаються й не додаються.

Результат чутливий до цього рішення. За основним правилом є 4035 внутрішніх і 1783 міжкласних ребра; за мітками контактів --- 4016 та 1802. Сума ваг міжкласних ребер змінюється з 0.674625 до 0.728258, тобто приблизно на 7.95\%. Відносна різниця матриць $\|W^{\rm contacts}-W^{\rm metadata}\|_F/\|W^{\rm metadata}\|_F$ дорівнює 25.41\%. Різниця сум і різниця матриць мають різні знаменники, тому ці відсотки не суперечать один одному. Це приклад того, чому навіть одна помилка належності до спільноти може змінити вхід макромоделі.

\paragraph{Часова повнота.}
Найраніший інтервал починається 2 грудня 2013 року о 11:00:00 UTC, останній закінчується 6 грудня о 15:59:40 UTC. Це межі опублікованих записів, а не гарантія безперервної роботи всіх сенсорів. Є 7375 інтервалів з хоча б одним контактом. Кількість таких інтервалів за днями: 899, 1619, 1619, 1619, 1619. Перший день має коротше покриття за цим показником. Відсутність рядків не дозволяє розрізнити відсутність контактів, ніч, відсутність людини та пропуск вимірювання. Тому підготовка не заповнює прогалини вигаданими контактами.

\section{Як часові контакти стають статичною мережею}
У поточній реалізації метод отримує статичні графи. Для кожної неорієнтованої пари визначаємо кількість контактних інтервалів $c_{uv}$ та сумарну тривалість
\[
d_{uv}=20c_{uv}\quad\text{секунд}.
\]
Вибираємо один спільний календарний знаменник для всіх пар:
\[
T_{\mathrm{obs}}=t_{\max}-t_{\min}+20=363\,580\ \text{с},\qquad
w_{uv}=\frac{d_{uv}}{T_{\mathrm{obs}}}.
\]
Доданок 20 потрібний тому, що мітки задають кінці інтервалів. Інакше втратили б довжину першого інтервалу. Ваги безрозмірні: це спостережена частка календарного часу з контактом. У знаменник входять ночі та інші прогалини; це обрана конвенція календарного усереднення.

Якщо $\beta$ має розмірність с$^{-1}$ за активного контакту, тоді в статичній моделі інтенсивність уздовж ребра дорівнює $\beta w_{uv}$ і вимірюється в с$^{-1}$. Наприклад, 300 секунд контакту дають $w=300/363580\approx0.0008251$. Довільно призначати цьому ребру вагу 300 і залишати ту саму $\beta$ означало б змінити модель у 363\,580 разів.

Варіанти $w'_{uv}=d_{uv}$ та $w''_{uv}=c_{uv}$ теж можливі, якщо відповідно взяти $\beta'=\beta/T_{\mathrm{obs}}$ та $\beta''=20\beta/T_{\mathrm{obs}}$. Це лише інші параметризації тих самих інтенсивностей. Натомість бінарна мережа $w^{\rm bin}_{uv}=1$ змінює відносні ваги сильних і слабких контактів і є іншою моделлю. Її можна завантажити окремо, але не слід змішувати з основною.

Для переходу на час у днях потрібно взяти $\beta_{\rm day}=86400\beta_{\rm sec}$ та $\gamma_{\rm day}=86400\gamma_{\rm sec}$; ваги залишаються тими самими. Крок синхронізації й горизонт також треба виразити в днях. Значення $\tau=1$ із синтетичних експериментів головної статті не стає автоматично ні однією секундою, ні одним днем у новому прикладі.

\paragraph{Що така агрегація зберігає.}
Для ребра вона зберігає сумарну контактну експозицію за повне вікно: $\beta w_{uv}T_{\mathrm{obs}}=\beta d_{uv}$. Це рівність для фіксованої пари, коли джерело весь час інфіковане, а одержувач сприйнятливий. В епідемії їхні стани змінюються всередині вікна, тому рівність загальної експозиції не доводить рівність розподілів зараження.

Втрачаються порядок зустрічей, паузи та сумісність контактів у часі. На статичному графі може існувати шлях, який неможливо пройти послідовно за реальним розкладом. Отже, майбутня перевірка двомасштабного методу тут має порівнювати його з повним CTMC на \emph{тій самій статичній зваженій мережі}. Порівняння з часовими даними є окремою задачею і додало б інше джерело похибки.

\section{Конкретна побудова двох рівнів}
Ідентифікатори людей залишаються оригінальними; перенумеровувати їх у $0,\ldots,N-1$ не потрібно. Класи впорядковуються лексикографічно, а ID усередині класу --- за зростанням. Це фіксує відповідність рядків $W$ локальним графам і запобігає прихованим перестановкам.

Спочатку будуємо неорієнтований граф усіх спостережених пар з вагами $w_{uv}$. Потім беремо індукований граф кожного класу. Кожне ребро між різними класами додає свою вагу до двох симетричних елементів $W$; внутрішні ребра в $W$ не потрапляють. Ребра не проріджуємо, класи не зрівнюємо за розміром і мережу не перетворюємо на ланцюг.

Отриманий граф має 327 вузлів і 5818 ребер та є зв'язним. Усі дев'ять локальних графів також зв'язні. Макрограф містить усі 36 можливих неорієнтованих міжкласних зв'язків; їхні ваги суттєво відрізняються. Тому аналог останнього класу $T_9$ із ланцюгового тесту тут не має природного змісту: немає одного наперед визначеного послідовного маршруту.

Виконуються перевірки
\[
\sum_iN_i=327,\quad \sum_i|E_{ii}|+\sum_{i<j}|E_{ij}|=4035+1783=5818,
\]
\[
\frac12\sum_{i,j}W_{ij}=\sum_{i<j}\sum_{\{u,v\}\in E_{ij}}w_{uv},\qquad
\sum_{\{u,v\}\in E}w_{uv}T_{\mathrm{obs}}=3\,770\,160\ \text{с}.
\]
Також перевірено симетрію $W$, нульову діагональ, додатність і скінченність усіх ваг. Це перевірки правильності підготовки, а не точності епідемічного наближення.

\section{Що можна зрозуміти зі структури до симуляцій}
\begin{table}[htbp]
\centering\small
\caption{Основний варіант: класи з метаданих, 327 контактних вузлів. ASPL обчислено без ваг; $\mu_2$ --- спектральний розрив зваженого нормованого лапласіана. $B_i$ містить людей із контактом хоча б в один інший клас.}
\begin{tabular}{lrrrrrr}
\toprule
Клас & $N_i$ & Внутр. ребра & ASPL & $\mu_2$ & $|B_i|$ & Внутр. сила, \%\\
\midrule
2BIO1 &36&402&1.373&0.1480&35&93.2\\
2BIO2 &34&385&1.317&0.0434&34&92.7\\
2BIO3 &40&593&1.240&0.1085&40&94.2\\
MP    &33&407&1.229&0.0823&33&90.1\\
MP*1  &29&211&1.488&0.0610&29&91.0\\
MP*2  &38&482&1.314&0.2020&38&93.4\\
PC    &44&678&1.283&0.2155&41&96.0\\
PC*   &39&534&1.279&0.1940&38&96.0\\
PSI*  &34&343&1.389&0.1943&33&90.4\\
\bottomrule
\end{tabular}
\end{table}

Загалом 93.49\% ваги контактів лежить усередині класів. У таблиці використано інший, локальний знаменник:
\[
f_i=\frac{2\sum_{\{u,v\}\in E_{ii}}w_{uv}}
{2\sum_{\{u,v\}\in E_{ii}}w_{uv}+\sum_{j\ne i}W_{ij}}.
\]
Множник 2 потрібний, бо внутрішнє ребро додає силу двом вузлам того самого класу, а зовнішнє --- одному. Значення $f_i$ показує частку сумарної сили вузлів класу, спрямовану всередину нього. Воно не є ймовірністю конкретного зараження.

\begin{figure}[htbp]
\centering
\includegraphics[width=0.83\linewidth]{figures/macro_contacts.pdf}
\caption{Матриця міжкласних контактів. Показано $W_{ij}T_{\mathrm{obs}}/60$, тобто сумарні парні хвилини, округлені до цілих. До симулятора передається безрозмірна $W$ з CSV. Діагональ порожня за означенням інтерфейсу: внутрішні контакти збережено в локальних графах. Дві симетричні клітинки представляють ту саму неорієнтовану суму; для загального підсумку треба брати один трикутник.}
\end{figure}

Дані підтримують вибір класів як змістовних спільнот. Але з переважання внутрішніх контактів не випливає швидке внутрішнє зараження. Наприклад, незважений ASPL класу 2BIO2 лише 1.317, проте його зважений $\mu_2=0.0434$ є найменшим серед класів. Короткі за кількістю ребер шляхи можуть складатися зі слабких контактів. Для обчислення $\mu_2$ ваги трактуються як сили зв'язків у
\[
L_{\rm norm}=I-D^{-1/2}AD^{-1/2},\qquad D_{uu}=\sum_v w_{uv}.
\]
Це структурний показник для зваженого випадкового блукання; він не дорівнює часу розповсюдження SI та не дає сам по собі оцінки похибки методу. Усі спектральні значення в таблиці отримані для зв'язних локальних графів.

\subsection*{Чому загальної кількості мостових вузлів недостатньо}
Майже всі люди мають хоча б один контакт поза власним класом. Однак важливо, \emph{з яким саме класом} вони зв'язані та яку частину ваги несуть. Для інтерфейсу $i\to j$ визначимо
\[
b^{(ij)}_u=\sum_{v\in V_j}w_{uv},\quad
p^{(ij)}_u=\frac{b^{(ij)}_u}{W_{ij}},\quad
N_{\rm eff}^{(ij)}=\frac{1}{\sum_{u\in V_i}(p^{(ij)}_u)^2}.
\]
За однакових ненульових ваг кінців $N_{\rm eff}$ дорівнює їх кількості; за концентрації на одній людині наближається до 1. Для сильного інтерфейсу MP--MP*2 є 130 ребер і відповідно 32 та 31 різні кінці, але ефективні числа становлять лише 7.25 та 2.82. Для 2BIO2--2BIO3 є 169 ребер, 31 та 36 кінців, а ефективні числа --- 12.95 та 4.58. Наявність багатьох мостів не означає рівномірного доступу до зовнішньої ваги.

Порівняємо всі міжкласні ваги з теоретичним рівномірним інтерфейсом:
\[
P_{uv}=\frac{w_{uv}}{W_{ij}},\qquad U_{uv}=\frac{1}{N_iN_j},\qquad
\operatorname{TV}_{ij}=\frac12\sum_{u\in V_i,v\in V_j}|P_{uv}-U_{uv}|.
\]
Для відсутнього ребра $w_{uv}=0$; такі пари обов'язково входять у суму. На 36 інтерфейсах TV лежить у межах 0.8863--0.9923, медіана 0.9767. Реальні ваги структурно далекі від повного рівномірного інтерфейсу. TV велика також через відсутні ребра, тому сама по собі не вимірює динамічну похибку.

Для фіксованого стану $x$ є проста межа. Множина активних пар $\{(u,v):x_u=I,x_v=S\}$ є подією в розподілах $P$ та $U$, звідси
\[
|a^\star_{ij}(x)-\lambda_{ij}(x)|\le
\beta_{\mathrm{out}}W_{ij}\operatorname{TV}_{ij}.
\]
Це абсолютна межа різниці миттєвих інтенсивностей; вона може бути грубою і не перетворюється автоматично на межу похибки часу приходу.

\begin{figure}[htbp]
\centering
\includegraphics[width=\linewidth]{figures/structure_checks.pdf}
\caption{Ліворуч: $100f_i$ для кожного класу. Праворуч: усі 36 неорієнтованих інтерфейсів, їхня кількість реальних ребер і TV до рівномірної сурогатної ваги. Це детерміновані характеристики одного графа; точки не є симуляціями та не мають статистичних довірчих інтервалів.}
\end{figure}

\section{Конкретний приклад різних інтенсивностей}
Розглянемо допустимий початковий стан: інфіковано лише учасника 763 із класу 2BIO2, усі інші сприйнятливі. У цього учасника немає контактних ребер до 2BIO3. Повна миттєва інтенсивність передачі від 2BIO2 до 2BIO3 дорівнює нулю. Натомість
\[
\frac{\lambda_{ij}}{\beta}=\frac{W_{ij}}{N_i}
=\frac{0.13532097475108643}{34}=0.0039800286691496>0.
\]
Це перевірено без запуску епідемії, за сусідами й формулою. Приклад вибрано перебором усіх одноінфікованих станів і пар класів із нульовою точною інтенсивністю так, щоб агрегована інтенсивність була найбільшою серед них. Це навмисний контрприклад, а не випадкова вибірка, оцінка частоти порушень чи доказ великої середньої часової похибки.

Він показує, яку інформацію втрачає $W$: для всіх людей класу макрошар бачить ту саму частку $1/N_i$, хоча їхні зовнішні сусіди різні. Для відновлення реальних маршрутів цього інтерфейсу недостатньо. Для розподілів часу приходу відповідь потребує симуляцій: локальна інфекція може швидко дістатися реальних мостів.

\section{Як співвіднести приклад із головною статтею}
Теоретична еквівалентність ідеального методу рівномірному сурогатному графу пояснює, що означає наша $W$. Натомість результат про стохастично ранній перший експорт у головній статті має припущення випадкового рівномірного розміщення мостів. Тут мости спостережені й неоднорідні, тому напрямок середнього зміщення з тієї теореми не переносимо.

У головній статті найменші розбіжності спостерігалися у 18 синтетичних конфігураціях з $\rho<0.25$, де нормована відстань Вассерштейна не перевищувала 4.5\%. Це результат для перевіреної сім'ї мереж і режимів. Для класів різного розміру, нерівних ваг, паралельних маршрутів та інших одиниць часу він не є гарантією.

Показник $\rho=\beta_{\mathrm{out}}\sigma_T$ потребує часів зараження вузлів. ASPL, $\mu_2$, $f_i$ чи $N_{\rm eff}$ не дають готового значення $\sigma_T$. Тут $\rho$ та похибки розподілів не обчислювалися. Для зваженого інтерфейсу характерний час $1/(\beta W_{ij})$ відрізняється між парами класів; одне значення $\beta\sigma_T$ може не описати всі інтерфейси однаково.

\paragraph{Чому просте зменшення $\beta$ не створить кращий режим.}
Для SI, якщо всі локальні та зовнішні інтенсивності помножити на константу $c$, повний та ідеальний агрегований процеси змінять масштаб часу: $T\mapsto T/c$. Відповідно $\beta\sigma_T$ збережеться. За однаково перемасштабованого горизонту відносна похибка не покращиться лише від спільного зменшення $\beta$. За фіксованого числового кроку може змінитися похибка синхронізації, а за фіксованого горизонту --- цензурування.

Якщо окремо досліджувати слабші зовнішні зв'язки, можна помножити всі міжкласні ваги повного графа \emph{і} $W$ на $r<1$, залишивши локальні ваги. Це визначить новий сценарій з тим самим розташуванням контактів, але іншою відносною інтенсивністю. Його слід позначити як змінену модель, а не видавати за вихідні вимірювання. Такі сценарії тут не створювалися й не запускалися.

\section{Що лежить у папці та як це завантажити}
Шляхи нижче відносні до \path{paper/network_preparation/}.
\begin{itemize}
\item \path{raw/}: два незмінені файли; \path{sources.json}: URL, умови та хеші.
\item \path{prepared/nodes.csv}: усі 329 ID, два варіанти класу та прапорець контактів.
\item \path{prepared/edges.csv}: 5818 пар, кількості інтервалів, секунди й ваги.
\item \path{prepared/full_graph.graphml}, \path{prepared/local_0.graphml}--\path{prepared/local_8.graphml}: повний і локальні графи основної версії.
\item \path{prepared/W.csv}, \path{prepared/community_nodes.json}: макроваги та порядок вузлів.
\item \path{prepared/W_binary.csv}: бінарна модель; \path{prepared/W_contact_labels.csv}: варіант належності учасника 28.
\item \path{prepared/communities.csv}, \path{prepared/interfaces.csv}, \path{prepared/summary.json}: перевірки й характеристики; \path{prepared/communities_contact_labels.csv}: альтернативний поділ.
\item \path{prepared/manifest.json}: версії бібліотек та хеші коду, входів і виходів; \path{preparation_runs.jsonl}: початок, статуси й завершення підготовки.
\item \path{figures/}: рисунки, маніфест і \path{analysis.json} з точним контрприкладом.
\end{itemize}

Відтворення з кореня репозиторію:
\begin{verbatim}
.venv/bin/python paper/network_preparation/prepare_network.py
\end{verbatim}
Команда читає локальні сирі файли. Якщо результати вже є та хеші збігаються, перевіряє їх і пропускає перебудову. Після зміни коду чи входу потрібен новий каталог через \texttt{--output}: старі результати не перезаписуються. Окремі команди для рисунків, тестів і PDF наведені в README.

Завантаження з Python у корені проєкту:
\begin{verbatim}
from paper.network_preparation.prepare_network import load_prepared

micro, full, W, comm_nodes = load_prepared()
m2, f2, W2, n2 = load_prepared(class_source="contacts")
m3, f3, W3, n3 = load_prepared(population="metadata")
m4, f4, W4, n4 = load_prepared(weight_mode="binary")
\end{verbatim}
Повертаються NetworkX та NumPy: \texttt{micro} містить внутрішні ребра, \texttt{full} --- всі, \texttt{comm\_nodes} --- узгоджені списки ID. API повного симулятора приймає \texttt{full\_graph=full}, \texttt{comm\_nodes=comm\_nodes}, \texttt{infection\_rate=beta}, \texttt{recovery\_rate=0}. API двомасштабного --- \texttt{W=W}, \texttt{micro\_graphs=micro}, \texttt{beta\_micro=beta}, \texttt{beta\_macro=beta}, \texttt{gamma=0}, а також явно обрані горизонт і крок. Це узгоджений SI на одних вагах; значення \texttt{beta} з контактів не оцінювали.

Базова перевірка потребує \texttt{alpha\_mode="off"}, \texttt{correction\_mode="off"}. Завантажувач задає нейтральні \texttt{alphas=1}; спектральні поправки не калібрувалися. Мости збережено в контрольному графі для аналізу, але вони не входять до формули базового макрошару. Саме завантаження не створює епідемічної траєкторії.

\section{Питання для наступної перевірки}
Підготовлена мережа є конкретним входом наявного API. Це ще не валідація точності двомасштабного процесу на ній. Для продовження потрібно зафіксувати $\beta$, одиницю часу, спільний горизонт і початкові умови, а потім запускати обидві моделі за одним протоколом. У цьому завданні симуляції не запускалися.

Порівнювати доцільно час першого приходу в кожний клас, число охоплених класів, ймовірність досягнення класу до горизонту та загальну динаміку. Початково інфікований клас слід позначати окремо: його час приходу нульовий за означенням. Через кілька маршрутів можна додатково розглядати час охоплення всіх класів; це інша величина, ніж прихід у наперед вибраний клас ланцюга.

Недосягнуті до горизонту класи залишаються цензурованими. Потрібно показувати їхню частку; показники лише серед завершених приходів описують умовну вибірку та можуть бути зміщені. Спільний seed у різних алгоритмах сам по собі не створює змістовно парні траєкторії. Для довірчих інтервалів потрібні повторення та описаний спосіб їх отримання.

Окремо варто перевіряти допустимість переходу $i\to j$: чи були в момент агрегованого експорту інфіковані реальні кінці саме цього інтерфейсу, та чи можливий контакт із фактичним імпортованим вузлом. Перша перевірка необхідна, але не достатня для повної допустимості. Обидві ознаки обчислюються на станах відповідної траєкторії; не можна брати час експорту одного запуску й час зараження мостів іншого.

Спочатку слід порівняти базові моделі на одному статичному графі, потім окремо перевірити крок синхронізації, мітку учасника 28 і варіант 329 вузлів. Моделювання вихідної часової мережі, калібрування під конкретну інфекцію та генерація відсутніх контактів є іншими задачами. Контактні записи дають структуру й експозицію, але не біологічну ймовірність передачі.


\end{document}
